\section{子环}

记号约定:$A$是子环.

\begin{definition}{子环}{}
	设$B\subseteq A$.若$B$是$A$的加法群的子群,也是$A^{\times}$的子幺半群,则称$B$是$A$的子群.
\end{definition}

\begin{remark}{}{}
	此时,$B$是环,典范嵌入$B\hookrightarrow A$是环同态.
\end{remark}

\begin{example}{$\Z$的子环}{}
	$\Z$的每个包含$1$的子群一定等于$\Z$,因此$\Z$的子环只有$\Z$和$\left\{0\right\}$.
\end{example}

\begin{example}{生成的子环}{}
	设$\left(A_{i}\right)_{i\in I}$是$A$的一族子环,则$\bigcap\limits_{i\in I}A_{i}$是$A$的子环.对$A$的子集$X$,称$\langle A\rangle\coloneq\bigcap\limits_{\substack{X\subseteq B\\ B\text{是子环}}}B$是$X$生成的子环.
\end{example}

\begin{example}{环的中心化子}{}
	对$A$的子集$X$,$C_{A}\left(X\right)$是$A$的子环.特别地,$Z\left(A\right)$是$A$的子环.
\end{example}

\begin{example}{$\Omega$-群的自同态环}{}
	设$\Omega$是集合,$G$是$\Omega$-群,$E=\End_{\mathsf{Grp}}\left(G\right)$,$F=\left\{f\in E\middle|f\text{是}\Omega\text{-等变映射}\right\}$,则$F$是$E$的子环.称$F$为该$\Omega$-群的自同态环.
	
	特别地,记$F_{1}$为$\left\{f_{\alpha}\middle|\alpha\in\Omega\right\}$生成的$E$的子群,则$F=C_{E}\left(F_{1}\right)$.
\end{example}



















